% Generated draft: requires mathematical and source review.
\documentclass[11pt]{article}
\usepackage[margin=25mm]{geometry}
\usepackage{fontspec}
\setmainfont{Latin Modern Roman}
\usepackage{amsmath,amssymb,mathtools,unicode-math}
\setmathfont{Latin Modern Math}
\usepackage{enumitem,needspace,xurl,hyperref}
\setlength{\parindent}{0pt}
\setlength{\parskip}{5pt}
\setlength{\emergencystretch}{3em}
\newcommand{\N}{\mathbb{N}}
\newcommand{\Z}{\mathbb{Z}}
\newcommand{\Q}{\mathbb{Q}}
\newcommand{\R}{\mathbb{R}}
\newcommand{\C}{\mathbb{C}}
\newcommand{\F}{\mathbb{F}}
\newcommand{\A}{\mathbb{A}}
\newcommand{\PP}{\mathbb{P}}
\newcommand{\E}{\mathbb{E}}
\newcommand{\eps}{\varepsilon}
\newcommand{\dd}{\,\mathrm{d}}
\begin{document}
\section*{UnsolvedMath ID: 800161}
\textbf{Source problem code:} catalogue entry
\subsection{Definitions and mathematical statement}
The source asks for an omega theorem for the quantity denoted by \(I_3(T)\), conditional on the ternary additive divisor conjecture. The supplied record does not define \(I_3(T)\), specify a version of that conjecture, or identify the coefficients of the main-term polynomial \(\mathcal P_9\). These omissions remain unresolved hypotheses and notation in this draft.

For a function \(E(T)\), interpret \(E(T)=\Omega(T^\Theta)\) in the magnitude sense:
\[
\exists c>0\;\exists (T_n)_{n\geq 1},\qquad
T_n>0,\quad T_n\longrightarrow+\infty,\quad
|E(T_n)|\geq cT_n^\Theta\quad\text{for every }n\geq1.
\]
This interpretation does not require errors of both signs; the supplied record does not specify a signed omega convention.

The requested statement is: assuming the ternary additive divisor conjecture intended by the source, prove that there exists a constant \(\Theta>0\) such that
\[
I_3(T)-T\mathcal P_9(\log T)=\Omega(T^\Theta)
\qquad(T\to+\infty),
\]
where \(\mathcal P_9\) is the fixed degree-nine main-term polynomial intended by the source. Equivalently, under the magnitude interpretation above, prove that there exist \(\Theta>0\), \(c>0\), and positive parameters \(T_n\to+\infty\) satisfying
\[
\left|I_3(T_n)-T_n\mathcal P_9(\log T_n)\right|
\geq cT_n^\Theta
\quad\text{for every }n\geq1.
\]
The polynomial is part of the intended main term, rather than a polynomial that may vary with \(T\) or with \(n\). A fully specified mathematical statement requires the source's definition and normalization of \(I_3\), its precise conjectural hypothesis, and its identification of \(\mathcal P_9\).
\subsection{Short English statement}
Assuming the intended ternary additive divisor conjecture, show that the difference between \(I_3(T)\) and its stated degree-nine logarithmic main term has magnitude at least a positive power of \(T\) along an unbounded sequence. The definitions of \(I_3\), the conjecture, and the main-term polynomial are missing from the supplied record.
\subsection{Sources}
\begin{itemize}
\item Local UnsolvedMath registry, permanent ID \(800161\), source problem code \texttt{catalogue entry}; title: \emph{Moments of L-functions --- Moments of the Riemann zeta function --- Problem 11}. The registry labels this record open; that label is supplied metadata, not an independently verified mathematical or literature assessment.
\item Hamieh, Kadiri, and Ng, \emph{Moments of L-functions Problem List} (2024), Problem 11, \url{https://arxiv.org/abs/2405.17800}. This is the source reference supplied by the registry; it has not been independently consulted for this draft.
\item The statement attributes the question to Nathan Ng, Lethbridge. The registry's proposer field lists Hamieh, Kadiri, and Ng, with proposed year 2024.
\item Supplied provenance: extraction method \texttt{arxiv\_latex}; access date 2026-10-16. This date records the supplied source-access provenance and does not represent independent source verification performed for this draft.
\end{itemize}
\end{document}
